\section{Melatih Model}

Setelah data latih (\textit{train}) disiapkan dan direduksi dimensinya menjadi empat fitur paling prediktif, tahap selanjutnya adalah melatih model \textit{machine learning}. Pada penelitian ini, pelatihan tidak menggunakan konfigurasi standar bawaan sistem (\textit{default}), melainkan menerapkan teknik optimasi parameter tingkat lanjut melalui \texttt{RandomizedSearchCV}. Teknik pencarian acak ini diinstruksikan untuk mencari dan menguji berbagai kombinasi parameter pengatur (\textit{hyperparameter tuning}) secara otomatis guna menemukan arsitektur model yang menghasilkan performa klasifikasi terbaik. Terdapat lima algoritma klasifikasi yang dilatih dan dioptimalkan secara komparatif, yakni sebagai berikut.

\begin{enumerate}
    \item \textbf{Decision Tree:} model klasifikasi dengan struktur pemecahan aturan secara berjenjang. Pencarian pengaturan optimal difokuskan pada batas maksimal tingkat kedalaman pohon aturan (\texttt{max\_depth}) dan jumlah minimal sampel observasi yang disyaratkan untuk membuat cabang aturan baru (\texttt{min\_samples\_split}). Algoritma ini juga disematkan nilai pembobotan penalti \texttt{class\_weight=\{0:1, 1:4\}} untuk memperbesar sensitivitas deteksi sistem terhadap penyewa kelas minoritas (label 1).
    \item \textbf{Random Forest:} algoritma gabungan (\textit{ensemble}) berbasis \textit{bagging} yang memproses banyak \textit{Decision Tree} secara paralel untuk meminimalisasi varians dan risiko \textit{overfitting}. Pencarian pengaturan optimal diuji pada variasi jumlah pohon yang dibangun (\texttt{n\_estimators}), tingkat kedalaman maksimal (\texttt{max\_depth}), batas minimal data pembelahan (\texttt{min\_samples\_split}), serta penempatan ujung daun (\texttt{min\_samples\_leaf}). Serupa dengan \textit{Decision Tree}, model ini menerapkan parameter pembobotan penalti dengan rasio 4:1.
    \item \textbf{K-Nearest Neighbors (KNN):} model berbasis spasial yang mengklasifikasikan data berdasarkan jarak kedekatan observasinya. Karena perhitungan metrik jarak sangat rentan terdistorsi oleh perbedaan rentang skala angka, fitur dinormalisasi terlebih dahulu menggunakan fungsi \texttt{StandardScaler} di dalam sebuah alur pemrosesan terintegrasi (\textit{pipeline}). Sistem kemudian dieksekusi untuk menguji jumlah tetangga terdekat (\texttt{n\_neighbors}) dan metode perhitungan bobot jarak (\texttt{weights}) yang paling optimal.
    \item \textbf{Support Vector Machine (SVM):} model geometris yang bertugas mencari \textit{hyperplane} (batas pemisah) terbaik antar kelas. Serupa dengan KNN, data wajib diseragamkan skalanya terlebih dahulu di dalam \textit{pipeline}, lalu dilatih menggunakan fungsi \textit{Radial Basis Function} (RBF) untuk menangkap pola sebaran data yang \textit{non-linear}. Sistem secara komputasional mencari nilai batas toleransi kesalahan (\texttt{C}) dan radius seberapa jauh pengaruh satu titik observasi tunggal terhadap batas klasifikasi (\texttt{gamma}) diiringi penerapan parameter \texttt{class\_weight} 4:1.
    \item \textbf{XGBoost:} algoritma gabungan berbasis \textit{boosting} tingkat lanjut yang membangun model secara sekuensial guna mengevaluasi dan memperbaiki rentetan kesalahan komputasi dari klasifikasi sebelumnya. Parameter ruang pencarian mencakup nilai kecepatan belajar model (\texttt{learning\_rate}), tingkat kedalaman maksimal (\texttt{max\_depth}), dan total tahapan pembentukan iterasi (\texttt{n\_estimators}). Untuk mengimbangi ketimpangan jumlah data operasional, model diatur secara eksplisit menggunakan parameter kontrol \texttt{scale\_pos\_weight=4}.
\end{enumerate}

Proses pelatihan seluruh algoritma tersebut dilakukan menggunakan himpunan data latih (\texttt{X\_train} dan \texttt{y\_train}). Selama proses pencarian parameter terbaik berlangsung, sistem menjalankan tahapan (\textit{cross-validation}) sebanyak tiga kali (\texttt{cv=3}) dan mengevaluasi kualitas setiap model menggunakan metrik \textit{F1-Macro}. Selain itu, penerapan strategi pembobotan penalti (\textit{Cost-Sensitive Learning}) pada semua algoritma memastikan bahwa proses perhitungan matematis tetap seimbang secara objektif. Langkah ini bertujuan untuk mencegah model klasifikasi menjadi bias akibat dominasi jumlah penyewa kelas normal (0) yang jauh lebih besar dibandingkan kelas melanggar (1).